% Figure from MATH 590 notes, section 27. Compile with the notes preamble.
\begin{tikzpicture}[scale=1.3, >={Stealth[length=4.5pt]}]
  % S_1: fixed throughout
  \fill[figB!12] (-1,0) circle (1);
  \draw[figB, thick] (-1,0) circle (1);
  \draw[figB!60, thin] (-2,0) arc (180:360:1 and 0.28);
  \draw[figB!60, thin, dashed] (0,0) arc (0:180:1 and 0.28);
  \node[font=\small, figB] at (-1.75,0.85) {$S_1$};
  \fill[black] (-2,0) circle (1.2pt) node[left, font=\scriptsize] {$q_1$};
  % S_2 minus q_2: collapses to p along great circles through q_2
  \fill[figR!5] (1,0) circle (1);
  \draw[black!60, thick] (1,0) circle (1);
  \node[font=\small, black!70] at (1.8,0.85) {$S_2$};
  \foreach \h in {0.9,0.45} {
    \draw[figR!80, thick, ->] (1.92,0) arc (0:150:0.92 and \h);
    \draw[figR!80, thick, ->] (1.92,0) arc (0:-150:0.92 and \h);
  }
  \draw[figR!80, thick, ->] (1.92,0) -- (0.2,0);
  \filldraw[fill=white, draw=figR, thick] (2,0) circle (1.7pt) node[right, font=\scriptsize, inner sep=3pt] {$q_2$};
  \fill[black] (0,0) circle (1.4pt);
  \node[font=\scriptsize, fill=white, inner sep=1pt] at (0,-1.25) {$p$};
  \draw[black!45, thin] (0,-1.13) -- (0,-0.08);
\end{tikzpicture}
